% Part 2 opening chapter (2 of 2): the tests of the virial coupling.
% Carried line by line from: The Virial Partition Across Wide Range of Physical Scales (25 Feb 2026; sections 4.4-4.6, 6.1-6.3, 7, 8, Tables 5-7, 9);
% Dark Matter and Dark Energy as Virial Partners (Mar 2026; sections 4.3, 4.4 Table 2, 7); Gravitational Decoherence and the Virial Partition
% (Mar 2026; section 5.3); Virial Efficiency and Effective Nonlinear Exponent (25 Feb 2026; section 7, Table 4).
% Corrections applied: errata V7, V8, V10, V16, V22, V24, V25, V26, V28, T13; VIRIAL_CHECK.md items 7, 8, 10, 16, 20, 22, 23, 26; ledger notes 10(f), 11(f), 11(h).
% Euclid stated only as in Chapter ch:latetime, Section sec:lt_euclid. Numbers: docs/verification/scripts/verify_virial_papers.py (sections G-J).

\chapter{Tests of the virial coupling}\label{ch:virial_tests}

\section*{Overview}
The coupling $\beta_m=\Omega_m/2$ predicts, with no free amplitude, $\mu_0=-0.136$ with $\Sigma_0=0$, a curved ramp in $f\sigma_8(z)$, a raised $E_G$,
two expansion rates ($H_0^{(\gamma)}=67.16$, $H_0^{(m)}=72.26$), and a growth-only $\Omega_m$ below the geometric one (Chapter~\ref{ch:virial}). These
predictions are on record before the measurements that will decide them. This chapter sets out what is measured now, what the next surveys measure, and what
would count against the coupling.

\section{The present record}\label{sec:vt_record}
Table~\ref{tab:vt_status} summarises the current observational status of the cosmological virial partition. All predictions are derived from the single
coupling $\beta_m=\Omega_m/2$ with no additional free parameters.

\begin{table}[htbp]\centering\small
\caption{Observational status of the cosmological virial partition. IAM values from the Level~2 chain and Eq.~\eqref{part2:eq:mu_virial} \calc; observed
values \observed.}\label{tab:vt_status}
\begin{tabularx}{\linewidth}{>{\raggedright\arraybackslash}X>{\raggedright\arraybackslash}p{2.6cm}>{\raggedright\arraybackslash}p{3.4cm}>{\raggedright\arraybackslash}X}
\hline
Observable & IAM prediction & Observed & Status\\\hline
$\sigma_8$ & $0.7998\pm0.0058$ & $0.802^{+0.022}_{-0.018}$ (KiDS-Legacy + DES Y3 + DESI + Pantheon+~\cite{Stolzner2025}) & $0.12\sigma$\\
$S_8$ & $0.822\pm0.011$ & $0.815^{+0.016}_{-0.021}$ (KiDS-Legacy~\cite{Wright2025}) & $0.33\sigma$; does not separate the models\\
$H_0$ (matter sector) & $72.26$ & $73.04\pm1.04$ (SH0ES~\cite{Riess2022}) & $0.75\sigma$\\
$H_0$ (photon sector) & $67.16\pm0.47$ & $67.36\pm0.54$ (Planck~\cite{Planck2018VI}) & $0.37\sigma$\\
$\Omega_m^{\rm growth}$ & $0.299$ at $z=0.5$ & no growth-only measurement yet & open\\
$f\sigma_8(z)$, six DESI DR1 bins & ramp, $-2.2$ to $-0.1\,\%$ & errors $8$--$13\,\%$ per bin~\cite{DESI2024V} & consistent; not yet separating\\
$\mu_0$ & $-0.136$, $\Sigma_0=0$ & $0.11^{+0.45}_{-0.54}$ (DESI FS+BAO+BBN~\cite{DESI2024VII}) and others & inside every interval (Chapter~\ref{ch:latetime})\\\hline
\end{tabularx}
\end{table}

\paragraph{Growth-rate data.} Both models are statistically consistent with the current redshift-space-distortion data; the IAM growth rate is within the
measurement uncertainties at all published redshifts. On the six DESI DR1 full-shape bins $\chi^2$ is $4.51$ ($\Lambda$CDM) and $5.24$ (IAM, MGCAMB form), and
on SDSS DR16 $6.19$ and $6.95$ (Chapter~\ref{ch:sectortension}). \calc\ Current precision is insufficient to distinguish the models on growth alone.

\paragraph{Sector census.} A census of published cosmological measurements classified by which sector they probe sorts them by the worldline rule
(Section~\ref{sec:vc_h0}). The matter ruler carries $f\sigma_8$ and redshift-space distortions, cluster masses and counts, peculiar velocities, Cepheids,
megamasers and the supernova ladder they calibrate; the photon ruler carries the CMB temperature and polarisation spectra, CMB lensing, Shapiro delay, VLBI,
the coherence of gravitational-wave signals, BAO angles and lensing time delays. The four independently named tensions in the literature --- the $S_8$
tension, the $H_0$ tension, the cluster hydrostatic mass bias, and the Planck cluster count deficit --- are all matter-sector anomalies consistent with a
single suppressed matter-sector coupling $\mu<1$ with the photon sector unmodified. \interp\ The probe-by-probe count by this rule is still to be made; the
lensing time delays, on the photon ruler, are the one photon-sector measurement that lies on the matter-sector value (Table~\ref{tab:vc_h0}). \openprob

\paragraph{The $\chi^2$ budget.} Across the $H_0$ probes the difference between IAM and a single-$H_0$ $\Lambda$CDM is the consequence of the sector
assignment (Section~\ref{sec:vc_h0}); across the growth data the two models are equal within the errors. A combined $\chi^2$ over domains therefore
measures the sector assignment, not the coupling, and is not used as evidence here.

\section{The $E_G$ statistic}\label{sec:vt_eg}
The $E_G$ statistic (Zhang et al.\ 2007~\cite{Zhang2007EG}; Reyes et al.\ 2010~\cite{Reyes2010}) is defined observationally as
\begin{equation}
E_G(z)\equiv\frac{\Upsilon_{gm}(k,z)}{\beta(z)\,\Upsilon_{gg}(k,z)},\label{eq:vt_eg}
\end{equation}
where $\Upsilon_{gm}$ is the galaxy--matter cross-power spectrum (from galaxy--galaxy lensing), $\Upsilon_{gg}$ the galaxy auto-power spectrum, and $\beta=f/b$
the redshift-space distortion parameter. For general relativity, $E_G=\Omega_m/f(z)$. For IAM --- which modifies only the growth rate, not the lensing
kernel --- the prediction is
\begin{equation}
E_G^{\rm IAM}(z)=\frac{\Omega_m}{f_{\rm IAM}(z)},\label{eq:vt_egiam}
\end{equation}
where $f_{\rm IAM}(z)$ is the growth rate from the IAM-modified perturbation equations. \derived\ This is the standard modified-gravity formula
$E_G=\Omega_m\Sigma/f$ with $\Sigma=1$: IAM's lensing kernel is unmodified, so the change enters only through the reduced $f_{\rm IAM}$. Since IAM
suppresses $f$ relative to general relativity, $E_G^{\rm IAM}(z)>E_G^{\rm GR}(z)$ at all redshifts: by $+3.6\,\%$ today, $+1.8\,\%$ at $z=0.3$ and $+1.1\,\%$ at
$z=0.5$ (Figure~\ref{fig:growth_vs_geometry}b). \prediction\ Published measurements (Reyes et al.\ 2010~\cite{Reyes2010}; Alam et al.\ 2017~\cite{Alam2017};
Singh et al.\ 2019~\cite{Singh2019}) carry errors of order 10\,\% or more, larger than the shift. A dedicated numerical comparison using
Eq.~\eqref{eq:vt_egiam} with the IAM growth curves is a priority near-term analysis; all required data are already published. \openprob

\section{The three-way cluster test}\label{sec:vt_cluster}
Define the cluster multi-observable ratios
\begin{equation}
R_1=Y_{\rm SZ}/M_{\rm lens},\qquad R_2=L_X/Y_{\rm SZ},\qquad R_3=L_X/M_{\rm lens}.\label{eq:vt_R}
\end{equation}
$Y_{\rm SZ}$ (thermal SZ pressure) and $L_X$ (X-ray hydrostatic equilibrium) probe the matter sector; $M_{\rm lens}$ (weak lensing) probes the photon sector. In
the Level~1 form, where $\mu$ multiplies the gravitational source, IAM predicts that $R_1$ and $R_3$ are suppressed by $\mu(z_{\rm cluster})$ relative to
$\Lambda$CDM, while $R_2$ is unaffected because it is the ratio of two matter-sector quantities (the $\mu$ factor cancels exactly). This pattern ---
suppression in matter-to-photon ratios, no suppression in matter-to-matter ratios --- is the signature of the dual-sector structure in that form.
\prediction\ In the Level~2 form, where the term enters the friction on perturbations and the Poisson equation is unmodified, the hydrostatic and lensing
masses of a relaxed cluster agree and the ratios are not suppressed; the test therefore separates the two implementations as well as the models
(Chapter~\ref{ch:threeway}). The eROSITA, Planck SZ, and DES weak lensing public catalogues already overlap sufficiently for this test; it requires no new
observations. \openprob

\section{Euclid}\label{sec:vt_euclid}
Euclid measures both sectors simultaneously from the same sky volume: the VIS camera (weak lensing, the photon sector) and the NISP spectrograph
(redshift-space distortions, the matter sector)~\cite{Euclid2025}. The predicted signature is $\mu/\Sigma=0.864/1.000$ today --- a 14\,\% matter-to-photon
discrepancy measured over 1.5 billion galaxies across 15,000\,deg$^2$. \prediction\ Among currently viable modified-gravity models this combination is
unusual: $f(R)$ gives $\mu>1$ with $\Sigma=1$, and the self-accelerating DGP branch, which gives $\mu<1$ with $\Sigma=1$, carries a ghost
(Chapter~\ref{ch:latetime}).

What Euclid can resolve depends on the scales it models and on the shape of $\mu(z)$, and is set out in Section~\ref{sec:lt_euclid}: Euclid's forecasts are
made for a template whose deviation lasts to higher redshift than IAM's, and for the IAM $\mu(z)$ the sensitivity runs from about $0.3\sigma$ with
conservative scale cuts to about $7\sigma$ with all probes taken to the smallest validated scales. A forecast made with the IAM $\mu(z)$ itself is the first
calculation still to be done. \openprob\ The clustering and weak-lensing products of the first data release are expected in mid-2027~\cite{ESA2026DR1}.

\section{DESI Year 5 and the shape of the ramp}\label{sec:vt_desi}
DESI Year 5 forecasts $f\sigma_8$ to about $1\,\%$ per $\Delta z=0.1$ bin at $0.8<z<1.3$ and $0.4$--$0.7\,\%$ aggregated over all redshifts~\cite{DESI2016}.
IAM predicts a specific curved suppression ramp, Eq.~\eqref{eq:vc_fs8}: $4.25\,\%$ at $z=0$, $2.19\,\%$ at $z=0.295$, $1.35\,\%$ at $z=0.5$, $0.41\,\%$ at $z=1$,
decreasing to $0.13\,\%$ at $z=1.491$ (linear growth, Planck 2018 background, same early amplitude), following the functional form of $E(a)=\exp(1-1/a)$.
\prediction\ This curved ramp is distinguishable at full DESI and Euclid precision from a constant rescaling, and the shape of the ramp is the activation
history of the kinetic half measured directly.

\section{The DESI phantom crossing}\label{sec:vt_phantom}
DESI DR2 BAO measurements, combined with the CMB and supernovae, prefer a dark-energy equation of state with $w_0>-1$ and $w_a<0$, with a phantom crossing at
$z_{\rm cross}\approx0.35$--$0.5$ depending on the supernova compilation~\cite{DESI2025}. \observed\ The preference comes from distances alone --- BAO,
supernova and CMB --- with no growth data in the fit, so it cannot come from fitting suppressed growth. Under IAM, a single-fluid $w_0w_a$CDM fitter applied to
data that sit on two rulers --- photon-sector BAO and CMB distances, and supernovae calibrated on the matter ruler --- can infer $w_{\rm eff}(z)\neq-1$ to
reconcile them. Whether the apparent phantom crossing is an artefact of applying a single-sector fitter to a two-sector universe is the two-ruler test of
Chapter~\ref{ch:sectortension}. \openprob

\paragraph{The LRG1 bin.} The DESI DR1 LRG1 bin at $z=0.51$ is the known outlier of the first release: its BAO distances drive much of the evolving
dark-energy preference, and independent analyses find that removing LRG1 brings the fits back toward $\Lambda$CDM concordance (Liu, Wang \& Zhao
2024~\cite{Liu2024LRG}; \'O Colg\'ain et al.\ 2025~\cite{Colgain2025}). \observed\ IAM changes growth, not distances, and its growth-only $\Omega_m$ at
$z=0.51$ is only 5.2\,\% below the geometric value; if LRG1 is confirmed by DR2 and Year~5 data at an anomalous value well beyond that, it would require an
additional explanation beyond the virial partition mechanism. \openprob

\section{Falsification criteria}\label{sec:vt_falsify}
IAM is falsified by any of the following outcomes:
\begin{enumerate}
\item Euclid measures $\mu_0=0$ at a precision that excludes the IAM $\mu(z)$, while the cosmological tensions disappear, indicating the tensions were
systematic rather than physical.
\item Euclid measures $\mu\neq1$ \emph{and} $\Sigma\neq1$ with correlated modifications, inconsistent with the specific $\mu<1$, $\Sigma=1$ dual-sector
structure.
\item Laboratory experiments detect gravitational decoherence of photons, directly contradicting $\Sigma=1$.
\item DESI Year 5 or Euclid $f\sigma_8(z)$ shows a constant (redshift-independent) suppression, rather than the curved ramp prescribed by
$E(a)=\exp(1-1/a)$. A constant rescaling fits many modified-gravity models; the specific functional form of $E(a)$ is the IAM prediction.
\item With $\Omega_m$ revised, growth data that still require the old $\beta_m$.
\item In the Level~1 form, the three-way cluster test finds $R_2=L_X/Y_{\rm SZ}$ suppressed, meaning the two matter-sector quantities do not cancel.
\end{enumerate}
\paragraph{What does not falsify IAM.} Individual numbers being off by modest amounts. A collapsed fraction or a virial efficiency assigned to a chosen halo
population is a description of that population; the product the theorem requires, $f_{\rm coll}\eta_{\rm vir}=1/2$, is not adjustable. The span of the
virial $1/2$ across 37 decades in scale cannot accidentally reverse.

\paragraph{The virial-partition reading.} The interpretation of the two halves inherits all of these criteria and adds the growth--geometry split of
Section~\ref{sec:vc_omega}. If Euclid measures $\mu_0$ consistent with zero at a precision that excludes the IAM $\mu(z)$, IAM is falsified and the
virial-partition interpretation does not apply at cosmological scales. If Euclid measures $\mu_0\approx-0.136$ with $\Sigma_0\approx0$, the combination
$\mu<1$, $\Sigma=1$ would constitute a detection of the sector split and strong evidence for the virial partition. \interp
\begin{itemize}
\item If $\mu_0\approx-0.136$ and $\Sigma_0\approx0$: virial partition confirmed at cosmological scales.
\item If $\mu_0\approx0$: IAM falsified; dark matter and dark energy are separate phenomena requiring independent explanations.
\item If the curved $f\sigma_8(z)$ ramp is confirmed: the kinetic half activation history is directly measured.
\end{itemize}

\section{Conclusions}\label{sec:vt_conclusions}
The virial theorem, verified for $1/r$ potentials from the hydrogen atom to the cosmic horizon across a wide range of physical scales, provides the
derivation chain for IAM's coupling constant: $\beta_m=\Omega_m/2$. This is a mathematical theorem, not a fit, not a numerological coincidence, and not a free
parameter. The physical interpretation --- that the virial $1/2$ partition divides gravitational energy between geometric entropy (spacetime curvature) and
informational entropy (gravitational decoherence) --- propagates from the quantum scale (Landauer's principle, $k_BT\ln2$ per bit) to the cosmic scale
($\mu_0=-0.136$ suppressing matter growth while leaving photons unmodified). \interp

Three observational comparisons stand. $\sigma_8=0.800$ agrees with the 2025 joint measurement at $0.1\sigma$. The $H_0$ sector split places the Cepheid
and CMB endpoints within $1\sigma$ of their sector-specific predictions, with the lensing time delays the open item. The growth data are consistent with
the coupling at their present precision. IAM's predictions --- $\mu_0=-0.136$, $\Sigma_0=0$, $\sigma_8=0.800$, $H_0^{(\gamma)}=67.16$,
$H_0^{(m)}=72.26$\,km\,s$^{-1}$\,Mpc$^{-1}$ --- are on record before the Euclid and DESI Year 5 releases that will decide them. \prediction
